% ================== СЕМЕСТРОВЫЙ ПРОЕКТ ==================
\documentclass[12pt,a4paper]{article}

\usepackage[utf8]{inputenc}
\usepackage[russian]{babel}
\usepackage{geometry}
\geometry{left=25mm,right=15mm,top=20mm,bottom=20mm}

\usepackage{setspace}
\onehalfspacing
\usepackage{indentfirst}
\setlength{\parindent}{1.25cm}

\usepackage{amsmath,amssymb}
\usepackage{graphicx}
\usepackage{float}
\usepackage{listings}
\usepackage{xcolor}
\usepackage{hyperref}
\usepackage{caption}
\usepackage{booktabs}

\graphicspath{{output/}}

\lstset{
  language=Python,
  basicstyle=\footnotesize\ttfamily,
  keywordstyle=\color{blue},
  commentstyle=\color{gray},
  stringstyle=\color{red},
  numbers=left,
  numberstyle=\tiny\color{gray},
  frame=single,
  breaklines=true,
}

\begin{document}

\thispagestyle{empty}
\begin{center}
МИНИСТЕРСТВО ОБРАЗОВАНИЯ РЕСПУБЛИКИ БЕЛАРУСЬ\\[0.5em]
Учреждение образования\\[0.5em]
\textbf{<<Белорусский государственный университет информатики и радиоэлектроники>>}\\[2em]
Факультет компьютерных систем и сетей\\[0.5em]
Кафедра электронных вычислительных средств
\end{center}

\vspace{2em}
\begin{center}
\textbf{ОТЧЁТ}\\[0.5em]
по семестровому проекту\\[0.5em]
по дисциплине\\[0.5em]
\textbf{<<Встраиваемые системы обработки изображений для задач мультимедиа>>}\\[1em]
\textbf{Тема: <<Система сжатия данных L2L на основе блочной лестничной параметризации ДКП--ОДКП>>}
\end{center}

\vspace{3em}
\hfill\begin{minipage}{0.55\textwidth}
Выполнил:\\
магистрант 1 курса магистратуры\\
факультета компьютерных систем и сетей\\
группы 556301\\[0.3em]
\textbf{Лебедевич Артём Владимирович}\\[1.2em]
Преподаватель:\\
доцент кафедры ЭВС, канд. техн. наук\\
\textbf{Ключеня Виталий Васильевич}
\end{minipage}

\vfill
\begin{center}
Минск, 2026~г.
\end{center}
\newpage

\tableofcontents
\newpage

\section{Цель работы}

Разработать программную и аппаратную модели системы трансформационного кодирования изображений по схеме \textbf{L2L} (lossless-to-lossy) на основе блочной лестничной структурной параметризации прямого и обратного дискретного косинусного преобразования (ДКП--ОДКП) с использованием арифметики с фиксированной запятой и алгоритма ДКП с минимальной вычислительной сложностью (Лофлер, 11 умножений). Продемонстрировать эффект <<шахматной доски>> и его устранение с помощью информационного SIB-блока.

\section{Теоретические сведения}

\subsection{Схема L2L}

Схема L2L объединяет режимы сжатия \textit{lossless} (без потерь) и \textit{lossy} (с контролируемыми артефактами) в едином трансформаторе. Требования: обратимость <<целое к целому>>, фиксированная запятая, перфективная реконструкция при наличии side information.

\subsection{Блочная лестничная параметризация}

Полифазная матрица банка фильтров представляется каскадом блоков лестничной структуры. Для ДКП--ОДКП используются параллельные диагональные преобразования:
\begin{equation}
\begin{bmatrix} y_0 \\ y_1 \end{bmatrix}
=
\mathrm{round}\left(
\begin{bmatrix} C & 0 \\ 0 & D \end{bmatrix}
\begin{bmatrix} x_0 \\ x_1 \end{bmatrix}
+
\begin{bmatrix} 0 \\ s \end{bmatrix}
\right),
\end{equation}
где $C$ --- матрица ДКП-II, $D$ --- матрица ОДКП-III ($D = C^\mathsf{T}$).

\subsection{Эффект <<шахматной доски>>}

Утечка постоянной составляющей (DC leakage) в фильтрах ОДКП приводит к потере регулярности 1-го порядка и появлению блочной сетки на реконструированном изображении. Для устранения используется \textbf{SIB-блок} (Side Information Block):
\begin{equation}
s_i = D \cdot s_{i-1}, \quad s_0 = 0.
\end{equation}
В режиме lossless передаётся $s_n$; в режиме lossy --- не передаётся.

\subsection{Алгоритм Лофлера}

Восьмиточечное ДКП-II реализуется за 11 умножений и 29 сложений (табл.~5.1 методички). Двумерное преобразование --- separable: строки, затем столбцы.

\section{Архитектура системы}

Проект размещён в каталоге \texttt{l2l\_dct/}. Основные модули:
\begin{itemize}
\item \texttt{python/dct/} --- ДКП/ОДКП (reference + Loeffler);
\item \texttt{python/ladder/} --- лестничная параметризация, SIB;
\item \texttt{python/l2l/} --- анализ, синтез, кодек;
\item \texttt{fpga/hls/} --- HLS-ядро для Pynq;
\item \texttt{matlab/} --- сверка fixed-point.
\end{itemize}

Формат fixed-point: Q16.16 для коэффициентов; узел округления ladder --- потеря младших 15 бит дробной части.

\section{Ход работы}

\subsection{Фаза 1: ДКП/ОДКП}

Реализованы матричный эталон и separable 2-D DCT/IDCT. Проверка: max ошибка Loeffler vs матрица $< 10^{-10}$ (float). Fixed-point round-trip 2-D: PSNR $> 55$~dB.

\subsection{Фаза 2: Эффект <<шахматной доски>>}

Изображение $512 \times 512$ разбивается на подполосы $\xi$ и $s$ (верхняя/нижняя половины). Для каждого блока $8 \times 8$:
\begin{itemize}
\item прямой путь: $y_i = \mathrm{C2D}(\xi)$, $t_i = \mathrm{D2D}(s)$ с округлением;
\item обратный путь без SIB: ошибки округления $\Rightarrow$ блочная сетка;
\item с SIB: side\_y, side\_t восстанавливают точные коэффициенты.
\end{itemize}

\begin{figure}[H]
\centering
\includegraphics[width=0.95\textwidth]{phase2_chessboard.png}
\caption{Реконструкция: исходное / без SIB / с SIB (аналог рис.~5.15 методички)}
\end{figure}

\subsection{Фаза 3--4: L2L кодек}

Lossless: кодирование $y_i$, $t_i$, SIB (zlib). Lossy: per-band квантование JPEG-матрицы, без SIB. Построен график rate--distortion.

\begin{figure}[H]
\centering
\includegraphics[width=0.7\textwidth]{phase4_rate_distortion.png}
\caption{Rate--distortion характеристика L2L lossy}
\end{figure}

\subsection{Фаза 5: FPGA (Pynq)}

HLS-модуль \texttt{loeffler\_dct\_8pt.cpp}: ap\_fixed$<$16,6$>$, pipeline II=1. Testbench сверяет результат с golden vector из Python. Драйвер Pynq: \texttt{l2l\_driver.py}, notebook \texttt{l2l\_overlay.ipynb}.

\section{Результаты}

\begin{table}[H]
\centering
\caption{Метрики на тестовом изображении $512 \times 512$}
\begin{tabular}{@{}lcc@{}}
\toprule
Режим & MSE & PSNR, dB \\
\midrule
Без SIB (analysis--synthesis) & $\approx 3.5 \times 10^{-2}$ & $\approx 62.7$ \\
С SIB (lossless) & $0$ & $\infty$ \\
2-D DCT round-trip (fixed) & $< 10^{-2}$ & $> 55$ \\
\bottomrule
\end{tabular}
\end{table}

На реконструированном изображении без SIB наблюдается характерная <<шахматная>> блочная структура; с SIB --- bit-exact восстановление.

\section{Выводы}

\begin{enumerate}
\item Реализована L2L-система сжатия на базе блочной лестничной параметризации ДКП--ОДКП с fixed-point арифметикой.
\item Подтверждено проявление эффекта <<шахматной доски>> при отсутствии SIB и его полное устранение при передаче side information.
\item Разработано HLS-ядро 8-точечного ДКП для Pynq и организована верификация по golden vectors.
\item Система поддерживает переключение lossless/lossy одним пайплайном.
\end{enumerate}

\section*{Приложение: запуск}

\begin{lstlisting}
cd l2l_dct
python3 run_phase2_chessboard.py
python3 tests/test_sib_chessboard.py
\end{lstlisting}

\end{document}
